\section{Timing, functional forms and calibration}\label{sec:qsp:setup}

Parts~\ref{part:sp} and~\ref{part:mkt} are stated in continuous time with general functional forms,
which is the right level of abstraction for the results that can be proved. This part restates the
planner's problem in discrete time with explicit functional forms, so that it can be solved
numerically; Part~\ref{part:qmkt} does the same for the market economy and adds the policy
experiments. Nothing in the economics changes. What changes is that every shadow price acquires a date
and that the accounting identities, which in continuous time are relations between flows at an instant,
become a simultaneous system inside the period. Section~\ref{sec:qsp:model:block} is about the latter,
which is the one genuinely new object in this part.

\subsection{Notation carried over and notation renamed}\label{sec:qsp:setup:notation}

Two symbols used in the earlier literature clash with symbols already spoken for here, and we rename
rather than overload. The capital share in the production function is written $\chi$, because $\alpha$
is the marginal recycling yield of~\eqref{eq:sp:setup:recycling-derivatives}. The curvature of the
felicity function is written $\eta$, matching $\eta\equiv-u''C/u'$ in~\eqref{eq:sp:sum:Cgrowth},
because $\sigma$ is the storage share of~\eqref{eq:sp:shorthand}. The exponents in the cost functions
are written $\xi^j$ rather than $\varepsilon^j$, because $\varepsilon\equiv a'x/a$ is the elasticity of
the recycling yield.

\subsection{Timing}\label{sec:qsp:setup:timing}

Each period $t=0,1,2,\dots$ runs as follows. The economy enters with the five stocks
$\left(K_t,S_t,X_t,P_t,M^K_t\right)$. Production, extraction, exploration, consumption, waste
collection, treatment and recycling all take place within the period, so that material recovered in
$t$ is available as input in $t$: this is what the continuous-time model says, and departing from it
would change the circularity multiplier rather than merely discretise it. The stocks are then updated
and the economy enters $t+1$.

Writing $G_t\equiv I_t+\delta K_t$ for gross capital formation, the laws of motion are
\begin{align}
K_{t+1} &= \left(1-\delta\right)K_t+G_t,
&
S_{t+1} &= S_t+D_t-N_t,
&
X_{t+1} &= X_t+D_t,
\label{eq:qsp:states-1}
\\
M^K_{t+1} &= \left(1-\delta\right)M^K_t+\Phi^I_tG_t,
&
P_{t+1} &= \left[1-\theta\!\left(P_t\right)\right]P_t+\Xi_t,
&&
\label{eq:qsp:states-2}
\end{align}
which are~\eqref{eq:sp:states-motion} discretised. Two timing conventions in~\eqref{eq:qsp:states-2}
deserve to be flagged because they show up later as discount factors on shadow prices. The material
released as waste in period $t$ is $\delta M^K_t$, the depreciation of the stock the economy
\emph{entered} the period with; the material embodied by this period's gross investment, $\Phi^I_tG_t$,
is not available for release until $t+1$. And emissions $\Xi_t$ add to the pollution stock only from
$t+1$ onwards, so a tonne emitted today does its damage from tomorrow.

\subsection{The discrete-time planner's problem}\label{sec:qsp:setup:problem}

\begin{problem}[The planner's problem in discrete time]\label{prob:qsp}
Choose $\left\{C_t,G_t,N_t,D_t,\varpi_t,K^R_t,W_t,\Omega_t\right\}_{t\ge0}$ to maximise
\begin{align}
U_0 &= \sum_{t=0}^{\infty}\beta^t\left[u\!\left(C_t\right)-v_t\!\left(P_t\right)\right],
&
0<\beta<1,
\label{eq:qsp:objective}
\end{align}
subject to, at every $t$, the goods constraint, the ledger and the intensity condition
\begin{align}
F_t\!\left(K_t-K^R_t,R_t,P_t\right) &= C_t+G_t+C^N_t+C^D_t+c^W_t\!\left(\varpi_t\right)W_t,
\label{eq:qsp:goods}
\\
W_t &= R_t-\Phi^I_tG_t+\delta M^K_t+\mathcal N_t,
\label{eq:qsp:ledger}
\\
R_t &= \Omega_t\left(\mathcal D_t+\phi^I_tG_t\right),
\qquad \Phi^I_t\equiv\Omega_t\phi^I_t,
\label{eq:qsp:intensity}
\end{align}
to the laws of motion~\eqref{eq:qsp:states-1}--\eqref{eq:qsp:states-2}, to the definitions
\begin{align}
R^R_t &= \mathcal R_t\!\left(\varpi_tW_t,K^R_t\right),
&
R_t &= N_t+R^R_t,
&
\Xi_t &= W_t\left(1-\varpi_t+d^W\varpi_t\right)-d^WR^R_t,
\label{eq:qsp:definitions}
\end{align}
with the waste bases $\mathcal D_t$ and $\mathcal N_t$ of~\eqref{eq:sp:setup:base}, and to the domain
restrictions $C_t,G_t,N_t,D_t,K^R_t,W_t,\Omega_t\ge0$ and $\varpi_t\in\left[0,1\right]$, given
$\left(K_0,S_0,X_0,P_0,M^K_0\right)$.
\end{problem}

The subscript $t$ on $F$, $\mathcal R$, $C^N$, $C^D$, $c^W$, $v$ and the intensity coefficients carries
the technical progress that the continuous-time statement suppressed. Everything else is
Problem~\ref{prob:sp} verbatim.

\subsection{Functional forms}\label{sec:qsp:setup:forms}

\smalltitle{Preferences.} Constant relative risk aversion in consumption, with a convex and
increasingly weighted disutility of pollution:
\begin{align}
u\!\left(C_t\right) &= \frac{C_t^{1-\eta}}{1-\eta},
&
v_t\!\left(P_t\right) &= \frac{\psi_t}{1+\nu}P_t^{1+\nu},
&
\psi_t &= \bar\psi\left(1+g^{\psi}\right)^t,
&
\eta,\nu,\bar\psi &>0,\ g^\psi\ge0 .
\label{eq:qsp:preferences}
\end{align}
The growth term $g^\psi>0$ says that the marginal willingness to pay for environmental quality rises
as conventional consumption grows. It is not innocuous: with $g^\psi=0$ the pollution term becomes
negligible relative to consumption along a growing path, and the model's environmental margins go
quiet in the long run for a reason that has nothing to do with the circular economy.

\smalltitle{Final goods.} Cobb--Douglas in capital and in materials above a physical floor, with
pollution as a productivity shifter:
\begin{align}
F_t &= A_t\left(K^Y_t\right)^{\chi}\left(R_t-\bar R\right)^{1-\chi},
&
A_t &= \bar A\left(1+g^A\right)^tP_t^{-\gamma},
&
0<\chi<1,\quad \gamma>0,\quad R_t>\bar R\ge0 .
\label{eq:qsp:production}
\end{align}
The floor $\bar R$ is the minimum material throughput the laws of physics impose on a given scale of
activity, the standard reply of \textcite{daly1977} and \textcite{stern1997} to unbounded
substitution between produced capital and natural inputs. Setting $\bar R=0$ recovers the ordinary
Cobb--Douglas case. The derivatives are $F_K=\chi Y/K^Y$, $F_R=\left(1-\chi\right)Y/\left(R-\bar R\right)$
and $F_P=-\gamma Y/P$.

\begin{remark}[What $\bar R$ does to the essentiality assumption]\label{rem:qsp:Rbar}
Proposition~\ref{prop:sp:phases} assumes $\lim_{R\to0}F_R=\infty$. With $\bar R>0$ the Inada point
moves: $F_R\to\infty$ as $R\downarrow\bar R$, so the proposition goes through with ``materials are
essential'' reading $R_t>\bar R$ at every date. The consequences are not cosmetic. First, the corner
$N=0$ of~\eqref{eq:sp:sum:N-corner} now requires the loop alone to deliver more than the floor,
$R^R_t>\bar R$, which is a much stronger requirement than $R^R_t>0$. Second, the long-run discussion of
Section~\ref{sec:sp:opt:results} acquires a hard constraint rather than an asymptotic one: extraction
cannot go to zero unless the circularity multiplier can hold $R$ above $\bar R$ against a vanishing
virgin flow. This is exactly the substitution limit the ecological-economics literature has in mind,
and in this model it bites through the materials loop rather than through the production function.
\end{remark}

\smalltitle{Recycling.} We use the yield function
\begin{align}
a_t\!\left(x\right) &= a^m_t\,\frac{x}{1+x},
&
a^m_t &= 1-\frac{b}{\left(1+g^R\right)^t},
&
b&>0,\quad g^R\ge0,
\label{eq:qsp:recycling}
\end{align}
so that $a^m_t$ is the recycling frontier at date $t$ and $g^R$ is the rate of technical progress that
moves it towards, but never to, complete recovery. The implied objects of
Section~\ref{sec:sp:setup:primitives} are all available in closed form:
\begin{align}
\varepsilon_t &= \frac{1}{1+x},
&
a'_t\!\left(x\right) &= \frac{a^m_t}{\left(1+x\right)^2},
&
\alpha_t &= a^m_t\left(\frac{x}{1+x}\right)^2 = \frac{a_t^2}{a^m_t},
&
a'_t\!\left(0^+\right) &= a^m_t .
\label{eq:qsp:recycling-derivatives}
\end{align}
Three properties are worth recording. First, the constant-returns form of the technology is a familiar
one: substituting $x=K^R/T$ into~\eqref{eq:qsp:recycling} gives
\begin{align}
\mathcal R_t\!\left(T,K^R\right) &= a^m_t\,\frac{TK^R}{T+K^R}
\;=\;a^m_t\left[T^{-1}+\left(K^R\right)^{-1}\right]^{-1},
\label{eq:qsp:recycling-crs}
\end{align}
a CES aggregator of treated waste and recycling capital with elasticity of substitution one half,
scaled by the frontier. Both inputs are essential, as Section~\ref{sec:sp:opt:phases} requires.
Second, $a'_t\!\left(0^+\right)=a^m_t<\infty$, which is the finiteness that
Proposition~\ref{prop:sp:phases}(iii) needs; the threshold~\eqref{eq:sp:phases:threshold} therefore
reads
\begin{align}
a^m_t\,V_t &\le F_{K,t}\left(1-\zeta_t\Omega^Y_t\right)
\qquad\Longleftrightarrow\qquad
K^R_t=0 ,
\label{eq:qsp:threshold}
\end{align}
as concrete a statement of the transition condition as one could want: the recycling frontier times the
virgin displacement value, against the marginal product of capital. Since $a^m_t$ rises with $g^R$
while $F_K$ falls with capital accumulation and $V_t$ rises with depletion, the inequality is
self-terminating on three counts at once. Third, the assumption $\lim_{x\to\infty}a=1$
of~\eqref{eq:sp:setup:recycling} holds only in the limit $t\to\infty$: at any date the frontier is
$a^m_t<1$, so the circularity multiplier is capped, exactly as flagged at the end of
Section~\ref{sec:sp:opt:results}. Remark~\ref{rem:qsp:tail} below draws out what that cap implies.

\begin{remark}[The recycling frontier bounds sustainable depletion]\label{rem:qsp:tail}
In a materially stationary economy with no displaced mass, \eqref{eq:sp:multiplier-identity} gives
$R_t=N_t/\left(1-a_t\varpi_t\right)$. Since $a_t\le a^m_t$ and $\varpi_t\le1$,
\begin{align}
R_t &\le \frac{N_t}{1-a^m_t}\;=\;\frac{N_t\left(1+g^R\right)^t}{b},
\label{eq:qsp:tail-bound}
\end{align}
so holding material throughput above the floor $\bar R$ requires
$N_t\ge b\bar R\left(1+g^R\right)^{-t}$: \emph{the virgin flow may decline at most at the rate of
technical progress in recycling}. The bound is an accounting consequence, not an optimality result, and
it is independent of the production function, of the elasticity of substitution between capital and
materials, and of how much capital is thrown at recycling. Under this parameterisation the long run is
therefore governed by $g^R$ rather than by capital deepening --- capital deepening runs into the
frontier $a^m_t$ and stops there --- which is a sharper version of the conjecture recorded in
Section~\ref{sec:sp:opt:results}, and one that the simulations can check directly.
\end{remark}

\smalltitle{Extraction and exploration.} Following the cost specifications carried in the paper draft,
\begin{align}
C^N_t &= \frac{c^N_t}{1+\xi^S}\,N_tS_t^{-\left(1+\xi^S\right)},
&
C^D_t &= \frac{c^D_t}{1+\xi^X}\,D_tX_t^{1+\xi^X},
&
\xi^S,\xi^X&>0,
\label{eq:qsp:mining-costs}
\\
c^N_t &= \frac{c^N}{\left(1+g^N\right)^t},
&
c^D_t &= \frac{c^D}{\left(1+g^D\right)^t},
&
g^N,g^D&\ge0 .
\label{eq:qsp:mining-progress}
\end{align}
Extraction gets harder as the reserve is drawn down, $C^N_S<0$; exploration gets harder as cumulative
discoveries accumulate, $C^D_X>0$; and cost-reducing technical progress shifts both schedules down.
Both are linear in the flow, so $C^N_{NN}=C^D_{DD}=0$, which satisfies the weak convexity assumed
in~\eqref{eq:sp:setup:extraction-cost}--\eqref{eq:sp:setup:discovery-cost} but leaves the flow margin
without curvature of its own.\footnote{This matters more for the numerics than for the economics. With
$C^N_{NN}=0$ the extraction condition~\eqref{eq:sp:sum:N} pins $N_t$ only through the dependence of
$\Psi_t$ on $R_t$, which can make the within-period solve of Section~\ref{sec:qsp:model:block}
ill-conditioned when the material demand schedule is flat. The obvious generalisation,
$C^N_t=\frac{c^N_t}{1+\xi^N}N_t^{1+\xi^N}S_t^{-\left(1+\xi^S\right)}$ with $\xi^N>0$, restores strict
convexity at the cost of one parameter, and we expect to need it.}

\smalltitle{Waste collection and treatment.} With $c^W_t\left(\varpi\right)=c^c_t+c^T_t\left(\varpi\right)$,
\begin{align}
c^c_t &= \frac{\bar c^c}{\left(1+g^c\right)^t},
&
c^T_t\!\left(\varpi\right) &= \frac{\bar c^T}{\left(1+\xi^T\right)\left(1+g^T\right)^t}\,\varpi^{1+\xi^T},
&
\xi^T&>0,\quad g^c,g^T\ge0 .
\label{eq:qsp:waste-costs}
\end{align}
This satisfies~\eqref{eq:sp:setup:waste-cost}: $c^T_t\left(0\right)=0$ and
$c^{T\prime}_t\left(\varpi\right)=\bar c^T\left(1+g^T\right)^{-t}\varpi^{\xi^T}$, which vanishes at
$\varpi=0$ and is increasing and convex. Note that
$c^{T\prime}_t\left(1\right)=\bar c^T\left(1+g^T\right)^{-t}$ is finite and \emph{falling} over time, so
the upper corner $\varpi_t=1$ flagged in~\eqref{eq:sp:sum:varpi-upper} is not only admissible under
this parameterisation but becomes progressively easier to reach. Full treatment is a case the solution
algorithm must handle, not an edge case to be assumed away.

\smalltitle{Material intensities.} The relative intensities and the durables coefficient decline with
technical and structural change,
\begin{align}
\omega^j_t &= \frac{\omega^j}{\left(1+g^j_\omega\right)^t},
\quad j\in\left\{Y,N,D,W,C\right\},
&
\phi^I_t &= \frac{\phi^I}{\left(1+g^\phi\right)^t},
&
g^j_\omega,g^\phi&\ge0 ,
\label{eq:qsp:intensities}
\end{align}
while the nature-attributable coefficients are tied to the state of the deposit,
\begin{align}
\Omega^{N,S}_t &= \bar\Omega^{N,S}\left(\frac{S_0}{S_t}\right)^{\xi^{G}},
&
\Omega^{D,S}_t &= \bar\Omega^{D,S},
&
\xi^{G}&\ge0 .
\label{eq:qsp:nature-intensities}
\end{align}
The exponent $\xi^G$ is the declining-ore-grade channel: as the reserve is drawn down, more overburden
and gangue must be moved per tonne of usable material. Setting $\xi^G=0$ switches it off, and the
baseline should probably do so until there is evidence to calibrate it against.

\begin{remark}[$\Omega$ is an outcome, and the intensities are identified only up to scale]\label{rem:qsp:scale}
Two points about~\eqref{eq:qsp:intensities}, both of which depart from how the paper draft treats these
coefficients. First, the average intensity $\Omega_t$ is \emph{not} specified here, because it is not a
parameter: it is determined residually by the intensity condition~\eqref{eq:qsp:intensity}, as
Section~\ref{sec:sp:setup:aggregatematerial} establishes. Only the relative intensities $\omega^j_t$ and
$\phi^I_t$ are primitives. Second, those primitives are identified only up to a common scale. Replacing
$\left(\omega^j_t,\phi^I_t\right)$ by $\left(\kappa\omega^j_t,\kappa\phi^I_t\right)$ for any $\kappa>0$
leaves $\Omega_t\mapsto\Omega_t/\kappa$ and leaves every object that enters an optimality condition ---
the absolute intensities $\Omega^j_t=\Omega_t\omega^j_t$, the durables intensity
$\Phi^I_t=\Omega_t\phi^I_t$, the storage share $\sigma_t$ and the waste flow $\Omega_t\mathcal D_t$ ---
exactly unchanged. The allocation is therefore invariant to the normalisation, and one coefficient must
be fixed: we set $\omega^C=1$, so that all relative intensities are read as waste per unit of final good
relative to consumption. What the data have to identify is the \emph{ratios} $\omega^j/\omega^C$ and
$\phi^I/\omega^C$, together with the growth rates.
\end{remark}

\smalltitle{Environmental absorption.} Finally,
\begin{align}
\theta\!\left(P_t\right) &= \bar\theta P_t^{-\xi^P},
&
0<\xi^P<1,
&&\text{so}&&
\hat\theta_t &\equiv \theta_t+\theta'_tP_t = \left(1-\xi^P\right)\theta_t ,
\label{eq:qsp:absorption}
\end{align}
where $\hat\theta_t$ is the marginal decay rate that Section~\ref{sec:sp:opt:foc} shows is the relevant
one. Since $\xi^P<1$ we have $0<\hat\theta_t<\theta_t$: absorption saturates, but never so far that the
effective discount rate on damages turns negative. This specification therefore rules out the tipping
point noted after~\eqref{eq:sp:costate:P} by construction, which is a modelling choice and should be
stated as one.

\subsection{Calibration}\label{sec:qsp:setup:calibration}

The calibration is not yet done and the data folder is empty; what follows records the targets and the
mapping from targets to parameters, so that the work can be organised against it.

\smalltitle{Initial condition.} Following the paper draft, the economy starts in an early linear stage
in which the environment absorbs the entire waste flow, so that the pollution stock is initially
stationary:
\begin{align}
P_1 = \left[1-\theta\!\left(P_0\right)\right]P_0+\Xi_0 = P_0
\qquad\Longleftrightarrow\qquad
\bar\theta P_0^{1-\xi^P} &= \Xi_0 .
\label{eq:qsp:calibration-P}
\end{align}
This is one equation linking $\bar\theta$, $P_0$ and $\xi^P$ to the initial emission flow, and it fixes
$\bar\theta$ given the other two. The point of starting here is that the transition to the circular
economy is then entirely an outcome: at $t=0$ the pollution price is small, the reserve is large,
capital is scarce, and~\eqref{eq:qsp:threshold} holds comfortably.

\smalltitle{Targets.} Grouped by the block they identify:
\begin{itemize}[leftmargin=1.6em,itemsep=0.15em]
\item \emph{Preferences and capital.} $\beta$, $\eta$, $\delta$, $\chi$ from the usual macro targets:
a real interest rate, an intertemporal elasticity of substitution, a depreciation rate and a capital
share. $\bar A$ normalises units.
\item \emph{Materials.} The ratios $\omega^j/\omega^C$ and $\phi^I/\omega^C$ from material-flow
accounts: waste generated per unit of final expenditure in each use, and the material content of gross
fixed capital formation. Given the ratios, $\Omega_0$ is an outcome, and its value is the model's
prediction for economy-wide domestic material consumption per unit of output --- a check, not a target.
\item \emph{Recycling and treatment.} $b$ and $g^R$ from an observed recycling rate together with a view
on the technical frontier; $\bar c^T$ and $\xi^T$ from the observed treatment share $\varpi_0$ and the
cost of raising it; $\bar c^c$ from the waste sector's share of value added.
\item \emph{Mining.} $c^N$, $\xi^S$ from extraction cost shares and the reserve-to-production ratio;
$c^D$, $\xi^X$ from exploration expenditure and the historical elasticity of discoveries to cumulative
discovery. $\bar\Omega^{N,S}$ from ore-grade data.
\item \emph{Environment.} $\bar\psi$, $\nu$ from a marginal damage estimate at $P_0$;
$g^\psi$ from an income elasticity of environmental willingness to pay; $\gamma$ from the output damage
of pollution; $\bar\theta$ from~\eqref{eq:qsp:calibration-P} and $\xi^P$ from evidence on absorption
saturation.
\item \emph{Substitution limit.} $\bar R$ has no obvious empirical counterpart and should be treated as
a scenario parameter rather than calibrated, varied from $\bar R=0$ upwards as
Remark~\ref{rem:qsp:Rbar} suggests.
\end{itemize}

\smalltitle{Growth rates.} The benchmark takes technical progress to be neutral across sectors,
\begin{align}
g^A=g^R=g^N=g^D=g^c=g^T=g^\phi=g^\psi=g>0,
&&
g^j_\omega=0 \ \ \text{for all }j,
\label{eq:qsp:neutral-progress}
\end{align}
so that no sector is favoured by assumption and the transition is driven by depletion and accumulation
alone. Departures from~\eqref{eq:qsp:neutral-progress}, one sector at a time, are the natural
comparative-dynamics exercise, and Remark~\ref{rem:qsp:tail} already says that $g^R$ should be expected
to matter more than the others.
